\documentclass[11pt]{article}
\usepackage[T1]{fontenc}
\usepackage[utf8]{inputenc}
\usepackage{amsmath,amssymb,amsthm,mathtools}
\usepackage[margin=1in]{geometry}
\usepackage[hidelinks]{hyperref}
\hypersetup{pdftitle={The v-Numbers of (x2,y2,zc,xyz) and Its Integral Closure for Every c at Least 2},pdfauthor={Xiaojian Zeng}}

\newtheorem{theorem}{Theorem}
\newtheorem{lemma}[theorem]{Lemma}
\newtheorem{corollary}[theorem]{Corollary}
\newcommand{\m}{\mathfrak m}
\newcommand{\NP}{\operatorname{NP}}
\newcommand{\Soc}{\operatorname{Soc}}
\newcommand{\Ass}{\operatorname{Ass}}

\title{The \(v\)-Numbers of \((x^2,y^2,z^c,xyz)\) and Its Integral Closure for Every \(c\geq 2\)}
\author{Xiaojian Zeng\\
PBC School of Finance, Tsinghua University, 100190\\
\texttt{xiaojianzeng@foxmail.com}}
\date{}

\begin{document}
\maketitle

\begin{abstract}
Let \(K\) be an arbitrary field and
\(I_c=(x^2,y^2,z^c,xyz)\subset K[x,y,z]\), where \(c\geq2\).
We prove the simultaneous formulas
\[
  v(I_c)=2,
  \qquad
  v(\overline{I_c})=\left\lceil\frac c2\right\rceil .
\]
The proof gives the weighted half-space description
\(\overline{I_c}=(x^ay^bz^d:c(a+b)+2d\geq2c)\) and classifies all
socle monomials of both quotients.  Consequently integral closure lowers
the \(v\)-number at \(c=2\), preserves it at \(c=3,4\), and raises it
exactly when \(c\geq5\).
\end{abstract}

\section{Statement and context}

Write \(R=K[x,y,z]\) with its standard grading.  For a homogeneous ideal
\(J\subset R\), its \(v\)-number is
\[
 v(J)=\min\bigl\{d:\text{there are }f\in R_d\text{ and }
 \mathfrak p\in\Ass_R(R/J)\text{ with }(J:f)=\mathfrak p\bigr\}.
\]
This invariant arose in the study of generalized minimum-distance functions
and socle degrees; see Cooper et al.~\cite{CooperEtAl2020}.  We use only the
displayed definition.  Standard background on monomial ideals and integral
closure may be found in \cite{HerzogHibi2011,HunekeSwanson2006}.

Biswas, Mandal and Phukan recently compared \(v(I)\) with
\(v(\overline I)\) for classes of monomial ideals
\cite{BiswasMandalPhukan2026}.  Their Theorem~3.2 treats ideals in two
variables, whereas their Example~3.3 considers
\[
  I_5=(x^2,y^2,z^5,xyz)
\]
and computes \(v(I_5)=2<v(\overline{I_5})=3\).  The result below determines
that entire one-parameter family, including the smaller endpoints.  This is
a comparison with the cited source, not a claim of priority beyond the
literature inspected here.

\begin{theorem}\label{thm:main}
For every field \(K\) and every integer \(c\geq2\),
\[
 v(I_c)=2,
 \qquad
 v(\overline{I_c})=\left\lceil\frac c2\right\rceil,
 \qquad
 I_c=(x^2,y^2,z^c,xyz)\subset K[x,y,z].
\]
More precisely, the complete lists of socle-monomial exponent triples are
\begin{align*}
 \operatorname{SocMon}(R/I_c)
   &=\{(1,1,0),(1,0,c-1),(0,1,c-1)\},\\
 \operatorname{SocMon}(R/\overline{I_c})
   &=\left\{(0,0,c-1),
       \left(1,0,\left\lceil\frac c2\right\rceil-1\right),
       \left(0,1,\left\lceil\frac c2\right\rceil-1\right)\right\}.
\end{align*}
\end{theorem}

\section{The integral closure as a weighted half-space}

For \((a,b,d)\in\mathbb R_{\geq0}^3\), put
\[
 w(a,b,d)=c(a+b)+2d,
 \qquad
 H_c=\{(a,b,d)\in\mathbb R_{\geq0}^3:w(a,b,d)\geq2c\}.
\]
The three pure-power exponent vectors lie on the boundary of \(H_c\).
They generate the whole half-space as a Newton polyhedron.  Indeed, if
\(u\in H_c\) and \(S=w(u)\), then \(t=2c/S\leq1\), the point \(tu\) is a
convex combination of \((2,0,0),(0,2,0),(0,0,c)\), and
\(u-tu\in\mathbb R_{\geq0}^3\).  Moreover,
\(w(1,1,1)=2c+2\), so the generator \(xyz\) does not enlarge this
polyhedron.  Thus
\[
 \NP(I_c)=H_c.\tag{1}
\]
The following direct argument proves the needed lattice statement over
every field, without taking (1) as a black box.

\begin{lemma}\label{lem:closure}
For every \(c\geq2\),
\[
 \overline{I_c}
 =\bigl(x^ay^bz^d:c(a+b)+2d\geq2c\bigr).\tag{2}
\]
\end{lemma}

\begin{proof}
For a nonzero polynomial \(f\), let \(\nu(f)\) be the least \(w\)-weight
of one of its monomial terms.  The weighted initial forms lie in the domain
\(K[x,y,z]\), so
\(\nu(fg)=\nu(f)+\nu(g)\).  Each generator of \(I_c\) has weight at least
\(2c\); hence every element of \(I_c^i\) has weight at least \(2ci\).

Suppose that \(f\) is integral over \(I_c\).  It satisfies
\[
 f^n+a_1f^{n-1}+\cdots+a_n=0,
 \qquad a_i\in I_c^i.
\]
If \(\nu(f)<2c\), then \(f^n\) has weight \(n\nu(f)\), while every other
summand has weight at least
\[
 2ci+(n-i)\nu(f)>n\nu(f).
\]
The nonzero lowest-weight initial form of \(f^n\) therefore cannot cancel,
a contradiction.  Thus every monomial term of \(f\) has weight at least
\(2c\), proving one inclusion in (2).

Conversely, let \(m=x^ay^bz^d\) have integer weight
\(S=c(a+b)+2d\geq2c\).  The product
\[
 (x^2)^{ca}(y^2)^{cb}(z^c)^{2d}
\]
uses exactly \(S\) generators and divides \(m^S\): the coordinate
differences are respectively
\(a(S-2c),b(S-2c),d(S-2c)\).  Hence \(m^S\in I_c^S\), and the monic
equation \(T^S-m^S=0\) shows that \(m\) is integral over \(I_c\).
The elements integral over an ideal form an ideal
\cite[Chapter~1]{HunekeSwanson2006}; consequently every polynomial whose
terms have weight at least \(2c\) belongs to \(\overline{I_c}\).
\end{proof}

\section{Socle classifications}

Both \(I_c\) and \(\overline{I_c}\) are proper \(\m=(x,y,z)\)-primary
ideals: each contains a power of every variable, and (2) excludes \(1\).
Thus \(\m\) is the only associated prime of either quotient.  It follows
directly from the definition that \(v(J)\) is the least total degree of a
nonzero class in
\[
 \Soc(R/J)=(J:\m)/J
\]
for either ideal \(J\).  Since membership in a monomial ideal is termwise,
the standard monomials form a basis of the quotient and its socle is spanned
by the standard monomials killed by \(x,y,z\).  In particular, a homogeneous
sum cannot produce a smaller degree through cancellation.

First consider \(I_c\).  A standard monomial \(u=x^ay^bz^d\) has
\(a,b\leq1\), \(d\leq c-1\), and not both \(a=b=1\) when \(d\geq1\).
For such a monomial, direct divisibility by the four generators gives
\begin{align*}
 xu\in I_c&\Longleftrightarrow a=1\ \text{or}\ (b=1\ \text{and}\ d\geq1),\\
 yu\in I_c&\Longleftrightarrow b=1\ \text{or}\ (a=1\ \text{and}\ d\geq1),\\
 zu\in I_c&\Longleftrightarrow d=c-1\ \text{or}\ a=b=1.\tag{3}
\end{align*}
Solving (3) is exhaustive.  The case \(a=b=0\) fails.  The cases
\((a,b)=(1,0),(0,1)\) force \(d=c-1\), while \(a=b=1\) forces \(d=0\)
by standardness.  Hence the first list in Theorem~\ref{thm:main} is complete.
Its total degrees are \(2,c,c\), and therefore
\[
 v(I_c)=2.\tag{4}
\]

Now consider \(\overline{I_c}\).  By Lemma~\ref{lem:closure}, a monomial
of weight \(W\) is standard exactly when \(W<2c\).  Multiplication by
\(x,y,z\) changes the weight by \(c,c,2\), respectively.  Thus the monomial
is in the socle exactly when
\[
 W<2c,\qquad W+c\geq2c,\qquad W+2\geq2c.
\]
Since \(2c-2\geq c\) for \(c\geq2\), these conditions reduce to
\[
 2c-2\leq W<2c.\tag{5}
\]
Set \(s=a+b\).  The upper inequality in (5) forces \(s\leq1\).  If
\(s=0\), the unique solution is \(d=c-1\).  If \(s=1\), then
\[
 c-2\leq2d<c,
\]
whose unique integer solution is
\(d=\lceil c/2\rceil-1\).  There are two choices for \((a,b)\), giving the
second list in Theorem~\ref{thm:main}.  Its total degrees are
\[
 c-1,\qquad \left\lceil\frac c2\right\rceil,
 \qquad \left\lceil\frac c2\right\rceil.
\]
Because \(c-1\geq\lceil c/2\rceil\) for \(c\geq2\),
\[
 v(\overline{I_c})=\left\lceil\frac c2\right\rceil.\tag{6}
\]
Equations (4) and (6), together with the two complete socle lists, prove
Theorem~\ref{thm:main}.

\begin{corollary}
One has
\[
 v(\overline{I_2})<v(I_2),
\]
equality for \(c=3,4\), and
\[
 v(I_c)<v(\overline{I_c})
 \quad\Longleftrightarrow\quad c\geq5.
\]
\end{corollary}

\begin{proof}
The ordered pair is \((2,1)\) at \(c=2\), is \((2,2)\) at \(c=3,4\),
and has second component at least \(3\) exactly when \(c\geq5\).
\end{proof}

\section{Finite check and assistance disclosure}

A deterministic exact-arithmetic script independently enumerated all
standard monomials and tested the three variable-colon conditions for every
\(2\leq c\leq20\).  It returned precisely the two symbolic lists in
Theorem~\ref{thm:main}.  This bounded calculation is a diagnostic, not a
proof of the universal quantifier.  The distinction between a proved
infinite class and an exact implementation aimed at further cases is also
made explicitly in the different reconstruction setting of
\cite{Zeng2026Multiplicity}; the quantified scope of a computation matters
in both settings.  Here the universal conclusion rests entirely on
Lemmas~\ref{lem:closure} and the symbolic solution of (3)--(5).

OpenAI Codex was used during proof discovery, literature triage, symbolic
case organization, and preparation of the \LaTeX{} manuscript.  Before
manuscript preparation, a separate fresh language-model context reconstructed
the frozen argument, attacked its endpoint and field-independence claims, and
reran the finite checker.  That audit is not human peer review.  No numerical
or language-model output is used as a premise in the proof above; all
universal claims follow from the displayed algebraic arguments.

\bibliographystyle{plain}
\bibliography{references}

\end{document}
